\documentclass[12pt,a4paper]{article}

\usepackage[T1]{fontenc}
\usepackage[polish]{babel}
\usepackage[a4paper,margin=2.5cm]{geometry}
\usepackage{graphicx}
\usepackage{float}
\usepackage{hyperref}
\usepackage{amssymb}
\usepackage{amsmath}
\usepackage{booktabs}
\usepackage{enumitem}
\usepackage{xcolor}
\usepackage{tikz}
\usetikzlibrary{arrows.meta,positioning}
\usepackage{pgfplots}
\pgfplotsset{compat=1.18}
\usepgfplotslibrary{groupplots,statistics}

\definecolor{colOX}{HTML}{1F77B4}
\definecolor{colPMX}{HTML}{FF7F0E}
\definecolor{colOXC}{HTML}{2CA02C}
\definecolor{colPMXC}{HTML}{D62728}

\pgfplotsset{
    wykres/.style={
        width=0.92\textwidth, height=6.2cm,
        tick label style={font=\small}, label style={font=\small},
        legend style={font=\small, draw=none, fill=none},
        grid=major, grid style={gray!18},
        every axis plot/.append style={line width=1pt},
    },
}

\newcommand{\coursename}{Wstęp do algorytmów ewolucyjnych}
\newcommand{\projecttitle}{Harmonogramowanie zadań}
\newcommand{\doctype}{Dokumentacja końcowa}
\newcommand{\authorone}{Wojciech Jasiński}
\newcommand{\authortwo}{Piotr Iśtok}
\newcommand{\projectdate}{Warszawa, \the\year}
\newcommand{\h}[1]{\par\medskip\noindent\textbf{#1}}

\begin{document}
\sloppy

\begin{titlepage}
    \centering
    \includegraphics[width=0.8\textwidth]{politechnika.png}\par
    \vspace{2.5cm}
    {\Large \coursename\par}\vspace{1.2cm}
    {\Huge \bfseries \projecttitle \par}\vspace{1cm}
    {\Large \itshape \doctype \par}\vspace{2cm}
    {\Large Autorzy:\par}\vspace{0.5cm}
    {\large \authorone\par}\vspace{0.3cm}
    {\large \authortwo\par}
    \vfill
    {\large \projectdate\par}
\end{titlepage}

% =====================================================================
\section{Wprowadzenie i opis problemu}
% =====================================================================
\textbf{Motywacja.} Celem projektu jest opracowanie i~implementacja algorytmu ewolucyjnego dla
problemu harmonogramowania zadań w~środowisku wieloprocesorowym z~ograniczeniem pierwszeństwa
(\textit{ang. Task Scheduling Problem with Precedence Constraints}). Zadanie należy do klasy
problemów NP-trudnych~\cite{ullman1975}, co uzasadnia wykorzystanie metaheurystyk.

\h{Definicja formalna.} Problem zdefiniowano jako czwórkę $(T, p, l, M)$, gdzie:
\begin{itemize}[itemsep=3pt, topsep=3pt]
    \item $T = \{t_1, t_2, \dots, t_n\}$ jest zbiorem $n$ zadań do wykonania;
    \item $p : T \times T \to \{0, 1\}$ koduje relację pierwszeństwa: $p(t_i, t_j) = 1$ oznacza, że zadanie $t_i$ musi zostać zakończone przed rozpoczęciem $t_j$. Zakładamy, że graf skierowany indukowany przez~$p$ jest acykliczny (DAG); w~przeciwnym razie zbiór rozwiązań dopuszczalnych jest pusty;
    \item $l: T \to \mathbb{N}$ przypisuje każdemu zadaniu czas trwania;
    \item $M \in \mathbb{N}$ jest stałą liczbą procesorów (grup), między które rozdzielamy zadania. Procesor wykonuje jedno zadanie naraz; nazwę „procesor'' stosujemy dla zgodności z~literaturą.
\end{itemize}
\noindent
\textbf{Funkcja celu.} Wyzwaniem jest taki podział zadań między $M$ procesorów, by zachować
wszystkie więzy pierwszeństwa i~zminimalizować czas zakończenia ostatniego zadania, czyli
\textit{makespan} ($C_{max}$). Harmonogram optymalny przy pełnym wykorzystaniu zasobów kończy proces
w~najkrótszym możliwym czasie. W~notacji Grahama problem zapisujemy jako $P\,|\,prec\,|\,C_{max}$.
Rysunek~\ref{fig:przyklad} pokazuje przykładowy graf zależności i~jego harmonogram.

\begin{figure}[H]
\centering
\begin{tikzpicture}[
    task/.style={circle, draw, minimum size=9mm, inner sep=0pt, font=\small, align=center},
    >={Stealth[length=2mm]}, edge/.style={->, draw=gray}, crit/.style={->, draw=colPMXC, very thick, dashed}]
  \node[task] (A) at (0,1.0)   {A\\[-2pt]\tiny 2};
  \node[task] (C) at (2.6,1.0) {C\\[-2pt]\tiny 3};
  \node[task] (B) at (0,-1.0)  {B\\[-2pt]\tiny 4};
  \node[task] (D) at (2.6,-1.0){D\\[-2pt]\tiny 4};
  \draw[->, draw=gray, thick] (A)--(C);
  \draw[crit] (B)--(C) node[midway, above left=-2pt, font=\tiny, text=colPMXC]{zasoby};
\end{tikzpicture}
\hspace{0.7cm}
\begin{tikzpicture}[font=\small]
  \def\gh{0.8}
  \draw[->] (0,-0.25) -- (7.6,-0.25) node[right]{$t$};
  \foreach \t in {0,...,7} \draw (\t,-0.2)--(\t,-0.3) node[below=-1pt]{\tiny\t};
  \node[left] at (0,1.0+\gh/2) {\footnotesize P1};
  \node[left] at (0,0.0+\gh/2) {\footnotesize P2};
  \draw[fill=colOX!20] (0,1.0) rectangle (2,1.0+\gh) node[midway]{A};
  \draw[fill=colOX!20] (2,1.0) rectangle (6,1.0+\gh) node[midway]{D};
  \draw[fill=colPMXC!30] (0,0.0) rectangle (4,0.0+\gh) node[midway]{\bfseries B};
  \draw[fill=colPMXC!30] (4,0.0) rectangle (7,0.0+\gh) node[midway]{\bfseries C};
  \draw[dashed, colPMXC] (7,-0.25) -- (7,2.0) node[above]{\footnotesize $C_{max}=7$};
\end{tikzpicture}
\caption{Przykład ($M=2$, liczby w~węzłach to czasy trwania). Po lewej graf zależności: szare,
ciągłe strzałki to więzy pierwszeństwa, a~czerwona, przerywana linia $B\!\to\!C$ to \emph{nie}
więz pierwszeństwa, lecz krawędź maszynowa (wąskie gardło zasobowe -- B i~C dzielą procesor P2).
Po prawej zdekodowany harmonogram dla permutacji \texttt{[A, B, D, C]}. Jedynym poprzednikiem C
w~relacji pierwszeństwa jest A, więc C jest technologicznie gotowe już w~$t=2$, jednak jego start
zostaje opóźniony do $t=4$ z~powodu zajętości procesora P2 przez zadanie B. O~całkowitym czasie
$C_{max}$ decyduje ścieżka maszynowa B$\to$C, dlatego CPM wyznaczane jest bezpośrednio na
harmonogramie.}
\label{fig:przyklad}
\end{figure}

% =====================================================================
\section{Metoda}
\label{sec:metoda}
% =====================================================================
\textbf{Ogólna koncepcja.} Zastosowano algorytm memetyczny~\cite{moscato1989}, który łączy
klasyczne dla algorytmów genetycznych globalne przeszukiwanie przestrzeni rozwiązań z~optymalizacją
lokalną opartą o~analizę ścieżki krytycznej. Podejście to ma zapobiegać powstawaniu wąskich gardeł
w~harmonogramach.

\h{Reprezentacja osobnika.} Osobnik koduje kolejność wykonywania zadań: chromosomem jest
permutacja o~długości $|T|$, zawierająca wszystkie zadania w~kolejności priorytetu umieszczania
w~harmonogramie. Reprezentacja nie zawiera bezpośredniego przydziału do procesorów -- wynika on
z~działania dekodera. Każda permutacja zgodna z~porządkiem topologicznym koduje rozwiązanie
dopuszczalne.

\h{Dekoder (\textit{list scheduling}).} Zamianę chromosomu na harmonogram i~wyznaczenie jego czasu
wykonuje zachłanny dekoder typu \textit{list scheduling}~\cite{graham1969}. Ponieważ procesory są
nierozróżnialne, dekoder przydziela kolejne gotowe zadania (zgodnie z~priorytetem z~chromosomu)
procesorom zwalniającym się najwcześniej, respektując relację pierwszeństwa~$p$. Wyliczone $C_{max}$
stanowi wartość funkcji przystosowania. Dekoder udostępnia też pełny harmonogram (czasy
rozpoczęcia/zakończenia oraz przydział do procesora), wykorzystywany przez moduł optymalizacji
lokalnej.

\h{Naprawa dopuszczalności (\textit{ready-list repair}).} Krzyżowanie i~mutacja mogą wygenerować
permutacje niezgodne z~porządkiem topologicznym. Mechanizm naprawy oparty o~listę zadań gotowych
przebudowuje chromosom do najbliższego porządku topologicznego, zachowując względną kolejność
z~oryginalnej permutacji. Dzięki temu cała populacja pozostaje dopuszczalna i~nie trzeba karać
rozwiązań niedopuszczalnych w~funkcji celu.

\h{Operatory ewolucyjne.} Część ewolucyjną zaimplementowano w~oparciu o~bibliotekę
DEAP~\cite{deap2012}.
\begin{itemize}[itemsep=3pt, topsep=3pt]
    \item \textbf{Selekcja:} turniejowa~\cite{miller1995}, zapobiegająca przedwczesnej dominacji pojedynczego osobnika.
    \item \textbf{Krzyżowanie:} operatory permutacyjne unikające dublowania zadań. Porównano krzyżowanie porządkowe (OX)~\cite{davis1985}, zachowujące względną kolejność, z~krzyżowaniem z~częściowym mapowaniem (PMX)~\cite{goldberg1985}, opartym o~mieszanie pozycyjne.
    \item \textbf{Mutacja:} operator zamiany (\textit{swap}) -- losowa zamiana miejscami dwóch zadań.
\end{itemize}

\h{Optymalizacja lokalna (CPM).} Kluczowym elementem metody memetycznej jest moduł oparty o~analizę
ścieżki krytycznej~\cite{kelley1959}. W~odróżnieniu od klasycznej CPM liczonej na samym grafie
zależności zastosowano \textbf{analizę świadomą harmonogramu}: ścieżkę krytyczną wyznaczamy na
faktycznie zdekodowanym harmonogramie, śledząc na każdym kroku poprzednika z~relacji pierwszeństwa
\emph{lub} z~zajętości procesora. Dzięki temu analiza wykrywa również wąskie gardła zasobowe
i~różni się między osobnikami. Zadania krytyczne otrzymują wyższy priorytet (operacja
\textit{pull-left}); jeśli poprawia to $C_{max}$, zmiana -- zgodnie z~ewolucją
Lamarckowską~\cite{whitley1994} -- zastępuje poprzednika w~populacji.

\h{Konfiguracje algorytmu.} Z~połączenia dwóch operatorów krzyżowania (OX/PMX) oraz obecności bądź
braku modułu CPM powstają cztery porównywane konfiguracje: \texttt{ox}, \texttt{pmx} (czysty AE)
oraz \texttt{ox\_cpm}, \texttt{pmx\_cpm} (pełny algorytm memetyczny).

% =====================================================================
\section{Plan eksperymentów}
\label{sec:eksperymenty}
% =====================================================================
\textbf{Cel i hipotezy badawcze.} Eksperymenty oceniają jakość rozwiązań algorytmu oraz weryfikują
dwie hipotezy:
\begin{itemize}[itemsep=1pt, topsep=2pt, leftmargin=3em]
    \item[(H1)] lokalna optymalizacja CPM istotnie obniża medianę $C_{max}$ względem czystego AE;
    \item[(H2)] krzyżowanie OX istotnie dominuje nad PMX na grafach rzadkich ($|E|/|T|^2 < 0{,}1$).
\end{itemize}

\h{Zestawy testowe.} Jako instancje posłużyły grafy ze zbioru \textit{Standard Task Graph Set}
(Kasahara Lab)~\cite{tobita2002} oraz grafy generowane syntetycznie, z~kontrolowaną liczbą zadań
i~gęstością. Dla małych instancji ($|T| \leq 30$) wyznaczono rozwiązanie optymalne solverem MILP
(CBC poprzez PuLP~\cite{mitchell2011}), co pozwala raportować lukę optymalizacyjną; dla większych odniesieniem jest dolne ograniczenie
$LB = \max(\lceil \text{całkowita praca}/M \rceil,\ \text{ścieżka krytyczna})$.

\h{Protokół statystyczny.} Dla każdej konfiguracji wykonano 15 niezależnych uruchomień (ziarna
$0$--$14$). Jest to nieco mniej, niż pierwotnie planowano -- ograniczył nas długi czas obliczeń --
jednak różnice pozostają istotne statystycznie. Pełna kampania (1200 uruchomień AE)
zajęła około 17 godzin czasu CPU przy zrównolegleniu, a~wąskim gardłem były warianty CPM na
największych i~najgęstszych instancjach. Miarą jakości jest najlepszy $C_{max}$
na uruchomienie; raportujemy medianę. Istotność różnic weryfikowano sparowanym testem kolejności
rangowanej Wilcoxona~\cite{wilcoxon1945} (implementacja \texttt{scipy}~\cite{virtanen2020}; te same
ziarna), z~korekcją Holma w~obrębie każdej rodziny hipotez, $\alpha = 0{,}05$. Punktem odniesienia są heurystyki listowe HLFET~\cite{adam1974}
i~MCP~\cite{wu1990} oraz czysty AE bez komponentu lokalnego.

\h{Zastrzeżenie dot. budżetu (H1).} Konfiguracje memetyczne wykonują dodatkowe ewaluacje funkcji
celu wewnątrz przeszukiwania CPM, dlatego porównanie prowadzimy przy równej liczbie \emph{pokoleń},
a~nie przy równym budżecie ewaluacji. Przewaga CPM jest więc częściowo skutkiem większego
efektywnego nakładu na pokolenie.

\h{Hiperparametry.} Parametry AE są globalne (jednakowe dla wszystkich instancji). Wartości
i~ich uzasadnienie -- zakresy ugruntowane w~literaturze algorytmów
genetycznych oraz prac o~harmonogramowaniu zadań na grafach zależności~\cite{dejong1990,
grefenstette1986, omara2010, hou1994} -- zestawia tabela~\ref{tab:hparams}. Ablacji hiperparametrów
nie przeprowadzono.

\begin{table}[H]
\centering\footnotesize\setlength{\tabcolsep}{6pt}
\begin{tabular}{llp{8.4cm}}
\toprule
parametr & wartość & podstawa \\
\midrule
prawd. krzyżowania $p_c$ & $0{,}8$ & klasyczne $0{,}6$~\cite{dejong1990}--$0{,}9$~\cite{grefenstette1986}; GA dla DAG $0{,}65$--$0{,}95$~\cite{omara2010, hou1994} \\
prawd. mutacji $p_m$ & $0{,}3$ & zamiana (\textit{swap}) na osobnika; GA permutacyjne $0{,}1$--$0{,}5$ \\
populacja & $50$ & \cite{dejong1990}; powszechny zakres $30$--$100$ \\
liczba pokoleń & $100$ & stały budżet (wspólna długość krzywych zbieżności) \\
rozmiar turnieju & $3$ & łagodny nacisk selekcyjny~\cite{miller1995} \\
próby CPM / osobnika & $20$ & budżet lokalnego przeszukiwania (warianty memetyczne) \\
niezależne ziarna & $15$ & ograniczone czasem obliczeń; poniżej konwencji $30$ \\
\bottomrule
\end{tabular}
\caption{Hiperparametry algorytmu i~ich uzasadnienie. Wartości są globalne dla wszystkich instancji.}
\label{tab:hparams}
\end{table}

% =====================================================================
\section{Wyniki}
\label{sec:wyniki}
% =====================================================================
\textbf{Jakość rozwiązań.} Tabela~\ref{tab:main} zbiera medianę $C_{max}$ dla czterech konfiguracji
wraz z~rozmiarem, gęstością i~odniesieniem. \texttt{pmx\_cpm} jest najlepsza lub równie dobra na
każdej instancji i~dorównuje lub pokonuje heurystyki HLFET/MCP na 19 z~20 (przegrywa tylko na
\texttt{robot\_t88}: mediana 697 vs 694, lecz najlepszy przebieg 671). Część instancji jest wymuszona -- każda konfiguracja
osiąga ograniczenie (np. \texttt{rand0014\_t500}: wszyscy $893$).

\begin{table}[H]
\centering\footnotesize\setlength{\tabcolsep}{4.5pt}
\resizebox{\textwidth}{!}{%
\begin{tabular}{lrrrrrrrrr}
\toprule
instancja & $|T|$ & gęstość & $M$ & ref. & HLFET/MCP & \texttt{ox} & \texttt{pmx} & \texttt{ox\_cpm} & \texttt{pmx\_cpm} \\
\midrule
\texttt{syn\_t15\_p010}       &  17 & 0{,}093 & 2 &  52 &   54 & \textbf{52}  & \textbf{52}  & \textbf{52}  & \textbf{52} \\
\texttt{syn\_t22\_p018}       &  24 & 0{,}090 & 3 &  44 &   48 & 46           & 45           & \textbf{44}  & \textbf{44} \\
\texttt{syn\_t30\_p030}       &  32 & 0{,}132 & 3 &  66 &   66 & \textbf{66}  & \textbf{66}  & \textbf{66}  & \textbf{66} \\
\texttt{syn\_t40\_p006}       &  42 & 0{,}044 & 4 &  52 &   55 & 55           & \textbf{52}  & \textbf{52}  & \textbf{52} \\
\texttt{syn\_t40\_p035}       &  42 & 0{,}157 & 4 & 108 &  108 & \textbf{108} & \textbf{108} & \textbf{108} & \textbf{108} \\
\texttt{syn\_t60\_p005}       &  62 & 0{,}033 & 4 &  72 &   74 & 81           & 73           & \textbf{72}  & \textbf{72} \\
\texttt{syn\_t60\_p015}       &  62 & 0{,}073 & 4 &  77 &   88 & 89           & 81           & 81           & \textbf{78} \\
\texttt{syn\_t100\_p004}      & 102 & 0{,}024 & 4 & 126 &  126 & 128          & \textbf{126} & 127          & \textbf{126} \\
\texttt{syn\_t100\_p012}      & 102 & 0{,}059 & 4 & 131 &  153 & 162          & 143          & 142          & \textbf{135} \\
\texttt{syn\_t150\_p004}      & 152 & 0{,}022 & 4 & 230 &  241 & 234          & \textbf{230} & 232          & \textbf{230} \\
\texttt{rand0000\_t50\_p008}  &  52 & 0{,}044 & 4 &  66 &   77 & 73           & 67           & \textbf{66}  & \textbf{66} \\
\texttt{rand0093\_t50\_p078}  &  52 & 0{,}360 & 4 &  69 &   73 & \textbf{73}  & \textbf{73}  & \textbf{73}  & \textbf{73} \\
\texttt{rand0063\_t100\_p002} & 102 & 0{,}017 & 4 & 136 &  137 & 137          & \textbf{136} & \textbf{136} & \textbf{136} \\
\texttt{rand0173\_t100\_p039} & 102 & 0{,}187 & 4 & 199 &  218 & 228          & 214          & 212          & \textbf{206} \\
\texttt{rand0071\_t300\_p001} & 302 & 0{,}006 & 4 & 423 &  428 & 425          & \textbf{423} & 424          & \textbf{423} \\
\texttt{rand0000\_t300\_p016} & 302 & 0{,}080 & 4 & 448 &  505 & 531          & 503          & 512          & \textbf{470} \\
\texttt{rand0014\_t500\_p020} & 502 & 0{,}098 & 8 & 893 &  893 & \textbf{893} & \textbf{893} & \textbf{893} & \textbf{893} \\
\texttt{robot\_t88\_p003}     &  90 & 0{,}017 & 4 & 621 &  694 & 814          & 743          & 729          & \textbf{697} \\
\texttt{sparse\_t96\_p001}    &  98 & 0{,}018 & 4 & 484 &  486 & 502          & 486          & 486          & \textbf{484} \\
\texttt{fpppp\_t334\_p002}    & 336 & 0{,}011 & 4 &1783 & 1903 & 2067         & 1896         & 2017         & \textbf{1843} \\
\bottomrule
\end{tabular}}
\caption{Mediana $C_{max}$ wg konfiguracji (pogrubienie -- najlepsza na instancję). Odniesienie:
optimum MILP dla $|T|\le 30$, w~pozostałych przypadkach dolne ograniczenie LB. Kolumna HLFET/MCP
podaje wynik heurystyk listowych -- obie pokrywają się na każdej instancji.}
\label{tab:main}
\end{table}

\textbf{Luka do dolnego ograniczenia (LB).} Rysunek~\ref{fig:gap} pokazuje lukę czterech konfiguracji do LB na
reprezentatywnych instancjach: CPM obniża słupek, PMX leży poniżej OX, a~\texttt{pmx\_cpm} jest
najniższy.

\begin{figure}[H]
\centering
\begin{tikzpicture}
\begin{axis}[
    wykres, ybar, bar width=4.2pt, ymin=0, ylabel={luka do LB [\%]},
    symbolic x coords={syn\_t22,syn\_t60\_p005,syn\_t100\_p012,rand0173,rand0000\_t300,robot\_t88,fpppp},
    xtick=data, x tick label style={rotate=35, anchor=east, font=\footnotesize},
    legend columns=4, legend style={at={(0.5,1.04)}, anchor=south}, enlarge x limits=0.08,
]
\addplot[fill=colOX,draw=colOX]     coordinates {(syn\_t22,4.55)(syn\_t60\_p005,12.5)(syn\_t100\_p012,23.66)(rand0173,14.57)(rand0000\_t300,18.53)(robot\_t88,31.08)(fpppp,15.93)};
\addplot[fill=colPMX,draw=colPMX]   coordinates {(syn\_t22,2.27)(syn\_t60\_p005,1.39)(syn\_t100\_p012,9.16)(rand0173,7.54)(rand0000\_t300,12.28)(robot\_t88,19.65)(fpppp,6.34)};
\addplot[fill=colOXC,draw=colOXC]   coordinates {(syn\_t22,0)(syn\_t60\_p005,0)(syn\_t100\_p012,8.4)(rand0173,6.53)(rand0000\_t300,14.29)(robot\_t88,17.39)(fpppp,13.12)};
\addplot[fill=colPMXC,draw=colPMXC] coordinates {(syn\_t22,0)(syn\_t60\_p005,0)(syn\_t100\_p012,3.05)(rand0173,3.52)(rand0000\_t300,4.91)(robot\_t88,12.24)(fpppp,3.37)};
\legend{ox,pmx,ox\_cpm,pmx\_cpm}
\end{axis}
\end{tikzpicture}
\caption{Luka mediany $C_{max}$ do LB (\%) dla czterech konfiguracji.}
\label{fig:gap}
\end{figure}

\textbf{H1 -- wpływ CPM.} CPM istotnie obniża medianę $C_{max}$ wszędzie tam, gdzie istnieje zapas:
\textbf{25/25 istotnych testów na korzyść CPM-on, żaden przeciw} ($p_{\text{Holm}}\le 0{,}05$). Efekt
rośnie z~trudnością instancji ($-85$ dla OX na \texttt{robot\_t88}). Zbieżność zbiorcza
(rys.~\ref{fig:convagg}) -- średnia luka do LB po wszystkich 20 instancjach -- pokazuje, że warianty
z~CPM zbiegają wyraźnie szybciej.

\begin{figure}[H]
\centering
\begin{tikzpicture}
\begin{axis}[wykres, height=5.6cm, xlabel={pokolenie}, ylabel={średnia luka do LB [\%]},
    xmin=0, xmax=100, ymin=0, legend pos=north east, legend style={font=\footnotesize}]
\addplot[colOX, thick]   table[x=gen, y=ox]      {data/conv_aggregate.dat};
\addplot[colPMX, thick]  table[x=gen, y=pmx]     {data/conv_aggregate.dat};
\addplot[colOXC, thick]  table[x=gen, y=ox_cpm]  {data/conv_aggregate.dat};
\addplot[colPMXC, thick] table[x=gen, y=pmx_cpm] {data/conv_aggregate.dat};
\legend{ox,pmx,ox\_cpm,pmx\_cpm}
\end{axis}
\end{tikzpicture}
\caption{Zbieżność zbiorcza: średnia luka do LB po wszystkich 20 instancjach. Warianty z~CPM
spadają najszybciej; \texttt{pmx\_cpm} kończy najniżej (1,7\%).}
\label{fig:convagg}
\end{figure}

\textbf{H2 -- OX vs PMX.} Hipoteza nie potwierdza się -- zachodzi relacja przeciwna. Na każdej
nieremisowej instancji PMX istotnie pokonuje OX: na 20 instancjach PMX wygrywa 15 razy, remisuje 5,
a~OX nie wygrywa nigdy; remisy występują tylko tam, gdzie oba operatory
osiągają ograniczenie. Marginesy bywają duże na trudnych grafach rzeczywistych (\texttt{fpppp}:
$-171$, \texttt{robot\_t88}: $-71$). Prawdopodobny powód: naprawa listy gotowych już wymusza
pierwszeństwo, więc atut OX w~zachowaniu kolejności jest stłumiony, a~pozycyjne mieszanie PMX
skuteczniej przeszukuje przestrzeń priorytetów. Różnice luki widać na rys.~\ref{fig:gap}, gdzie
słupek PMX leży poniżej słupka OX na każdej instancji.

\textbf{Rozrzut.} Boxploty (rys.~\ref{fig:box}) pokazują rozrzut po ziarnach: \texttt{pmx\_cpm} jest
nie tylko najniższy, ale i~najwęższy.

\begin{figure}[H]
\centering
\begin{tikzpicture}
\begin{groupplot}[
    group style={group size=2 by 2, vertical sep=1.7cm, horizontal sep=1.5cm},
    width=0.5\textwidth, height=4.4cm, boxplot/draw direction=y,
    xtick={1,2,3,4}, xticklabels={ox,pmx,ox\_cpm,pmx\_cpm},
    x tick label style={font=\footnotesize, rotate=25, anchor=east},
    tick label style={font=\footnotesize}, ylabel style={font=\small}, ylabel={$C_{max}$},
    ymajorgrids=true, grid style={gray!18},
]
\nextgroupplot[title={\small \texttt{syn\_t100\_p012}}]
\addplot[draw=colOX,fill=colOX!25,boxplot prepared={lower whisker=158,lower quartile=161,median=162,upper quartile=163,upper whisker=164}]coordinates{};
\addplot[draw=colPMX,fill=colPMX!25,boxplot prepared={lower whisker=139,lower quartile=142,median=143,upper quartile=143.5,upper whisker=147}]coordinates{};
\addplot[draw=colOXC,fill=colOXC!25,boxplot prepared={lower whisker=138,lower quartile=141.5,median=142,upper quartile=142.5,upper whisker=144}]coordinates{};
\addplot[draw=colPMXC,fill=colPMXC!25,boxplot prepared={lower whisker=135,lower quartile=135,median=135,upper quartile=135,upper whisker=135}]coordinates{};
\nextgroupplot[title={\small \texttt{rand0173\_t100\_p039}}]
\addplot[draw=colOX,fill=colOX!25,boxplot prepared={lower whisker=224,lower quartile=227,median=228,upper quartile=230,upper whisker=231}]coordinates{};
\addplot[draw=colPMX,fill=colPMX!25,boxplot prepared={lower whisker=211,lower quartile=212.5,median=214,upper quartile=216,upper whisker=217}]coordinates{};
\addplot[draw=colOXC,fill=colOXC!25,boxplot prepared={lower whisker=209,lower quartile=211,median=212,upper quartile=212,upper whisker=215}]coordinates{};
\addplot[draw=colPMXC,fill=colPMXC!25,boxplot prepared={lower whisker=206,lower quartile=206,median=206,upper quartile=207,upper whisker=208}]coordinates{};
\nextgroupplot[title={\small \texttt{robot\_t88\_p003}}]
\addplot[draw=colOX,fill=colOX!25,boxplot prepared={lower whisker=784,lower quartile=810.5,median=814,upper quartile=823.5,upper whisker=834}]coordinates{};
\addplot[draw=colPMX,fill=colPMX!25,boxplot prepared={lower whisker=723,lower quartile=738,median=743,upper quartile=754.5,upper whisker=768}]coordinates{};
\addplot[draw=colOXC,fill=colOXC!25,boxplot prepared={lower whisker=711,lower quartile=723.5,median=729,upper quartile=733.5,upper whisker=736}]coordinates{};
\addplot[draw=colPMXC,fill=colPMXC!25,boxplot prepared={lower whisker=671,lower quartile=691.5,median=697,upper quartile=698.5,upper whisker=701}]coordinates{};
\nextgroupplot[title={\small \texttt{fpppp\_t334\_p002}}]
\addplot[draw=colOX,fill=colOX!25,boxplot prepared={lower whisker=2038,lower quartile=2060,median=2067,upper quartile=2079.5,upper whisker=2090}]coordinates{};
\addplot[draw=colPMX,fill=colPMX!25,boxplot prepared={lower whisker=1836,lower quartile=1870,median=1896,upper quartile=1930.5,upper whisker=1974}]coordinates{};
\addplot[draw=colOXC,fill=colOXC!25,boxplot prepared={lower whisker=1976,lower quartile=2008.5,median=2017,upper quartile=2026,upper whisker=2053}]coordinates{};
\addplot[draw=colPMXC,fill=colPMXC!25,boxplot prepared={lower whisker=1801,lower quartile=1821,median=1843,upper quartile=1869,upper whisker=1889}]coordinates{};
\end{groupplot}
\end{tikzpicture}
\caption{Rozrzut $C_{max}$ po 15 ziarnach. \texttt{pmx\_cpm} -- najniższy i~najwęższy.}
\label{fig:box}
\end{figure}

\textbf{Ujęcie zbiorcze.} Rysunek~\ref{fig:agg} porządkuje konfiguracje średnią luką do LB --
osobno dla wszystkich instancji, grafów gęstych i~rzadkich. W~każdym podziale CPM obniża lukę, PMX
leży poniżej OX, a~\texttt{pmx\_cpm} jest najlepszy. Efekt CPM jest wyraźnie silniejszy na grafach
rzadkich (OX $9{,}2\%$ wobec \texttt{pmx\_cpm} $1{,}6\%$) niż na gęstych, gdzie część instancji jest
wymuszona i~różnice się spłaszczają.

\begin{figure}[H]
\centering
\begin{tikzpicture}
\begin{groupplot}[group style={group size=3 by 1, horizontal sep=1.0cm},
    width=0.32\textwidth, height=5.2cm, ybar, ymin=0, ymax=10,
    symbolic x coords={ox,pmx,ox\_cpm,pmx\_cpm}, xtick={ox,pmx,ox\_cpm,pmx\_cpm}, enlarge x limits=0.18,
    x tick label style={rotate=25, anchor=east, font=\footnotesize}, tick label style={font=\footnotesize},
    title style={font=\small}, label style={font=\small},
    ymajorgrids=true, grid style={gray!18},
    nodes near coords, every node near coord/.append style={font=\tiny},
    every axis plot/.append style={bar shift=0pt, bar width=9pt}]
\nextgroupplot[title={wszystkie (20)}, ylabel={średnia luka do LB [\%]}]
\addplot[fill=colOX,draw=colOX]     coordinates {(ox,8.34)};
\addplot[fill=colPMX,draw=colPMX]   coordinates {(pmx,3.58)};
\addplot[fill=colOXC,draw=colOXC]   coordinates {(ox\_cpm,3.65)};
\addplot[fill=colPMXC,draw=colPMXC] coordinates {(pmx\_cpm,1.71)};
\nextgroupplot[title={gęste ($\ge 0{,}1$)}]
\addplot[fill=colOX,draw=colOX]     coordinates {(ox,5.09)};
\addplot[fill=colPMX,draw=colPMX]   coordinates {(pmx,3.33)};
\addplot[fill=colOXC,draw=colOXC]   coordinates {(ox\_cpm,3.08)};
\addplot[fill=colPMXC,draw=colPMXC] coordinates {(pmx\_cpm,2.33)};
\nextgroupplot[title={rzadkie ($< 0{,}1$)}]
\addplot[fill=colOX,draw=colOX]     coordinates {(ox,9.15)};
\addplot[fill=colPMX,draw=colPMX]   coordinates {(pmx,3.64)};
\addplot[fill=colOXC,draw=colOXC]   coordinates {(ox\_cpm,3.79)};
\addplot[fill=colPMXC,draw=colPMXC] coordinates {(pmx\_cpm,1.55)};
\end{groupplot}
\end{tikzpicture}
\caption{Średnia luka do LB wg konfiguracji: wszystkie instancje, grafy gęste i~rzadkie. CPM obniża
lukę w~każdym podziale; przewaga CPM jest największa na grafach rzadkich.}
\label{fig:agg}
\end{figure}

\textbf{Koszt obliczeniowy.} Rysunek~\ref{fig:cpuscatter} przedstawia średni czas CPU uruchomienia
w~funkcji rozmiaru instancji (skala log--log). Czysty AE i~warianty memetyczne tworzą dwie
równoległe chmury: koszt obu rośnie potęgowo z~rozmiarem w~zbliżonym tempie, lecz CPM dokłada do
każdego pokolenia powtórne dekodowanie harmonogramu, co przekłada się na w~przybliżeniu \emph{stały
współczynnik} (krotność) kosztu -- równoległe proste w~skali log--log -- stąd warianty memetyczne są
o~rzędy wielkości droższe na całym zakresie.

\begin{figure}[H]
\centering
\begin{tikzpicture}
\begin{axis}[
    width=0.92\textwidth, height=7.2cm, xmode=log, ymode=log,
    xlabel={liczba zadań $|T|$}, ylabel={średni czas CPU / uruchomienie [s]},
    tick label style={font=\small}, label style={font=\small},
    grid=major, grid style={gray!18},
    xmin=14, xmax=560, ymin=0.05, ymax=2000,
    legend pos=south east, legend style={font=\footnotesize, fill=white},
    legend cell align=left,
]
% --- power-law fits (linear regression in log-log space) ---
\addplot[colOX, thick, domain=15:520, samples=40, forget plot]    {0.001307*x^1.522};
\addplot[colPMX, thick, domain=15:520, samples=40, forget plot]   {0.001803*x^1.417};
\addplot[colOXC, thick, domain=15:520, samples=40, forget plot]   {0.09674*x^1.358};
\addplot[colPMXC, thick, domain=15:520, samples=40, forget plot]  {0.1184*x^1.318};
% --- measured points ---
\addplot[only marks, mark=*, colOX, mark size=1.7pt]            table[x=tasks,y=ox]      {data/cpu_time.dat};
\addplot[only marks, mark=square*, colPMX, mark size=1.6pt]     table[x=tasks,y=pmx]     {data/cpu_time.dat};
\addplot[only marks, mark=triangle*, colOXC, mark size=2pt]     table[x=tasks,y=ox_cpm]  {data/cpu_time.dat};
\addplot[only marks, mark=diamond*, colPMXC, mark size=2pt]     table[x=tasks,y=pmx_cpm] {data/cpu_time.dat};
\legend{ox,pmx,ox\_cpm,pmx\_cpm}
\end{axis}
\end{tikzpicture}
\caption{Czas CPU vs liczba zadań (skala log--log). Markery -- pomiary, linie -- dopasowanie
potęgowe $t=a\,|T|^{b}$.}
\label{fig:cpuscatter}
\end{figure}

% =====================================================================
\section{Dyskusja i wnioski}
\label{sec:wnioski}
% =====================================================================
\begin{itemize}[itemsep=3pt, topsep=3pt]
    \item \textbf{Najlepsza metoda to \texttt{pmx\_cpm}.} Nie wystarczy stwierdzić, że AE z~CPM jest
    dobry. Najlepsze wyniki daje konkretne połączenie PMX z~lokalnym przeszukiwaniem CPM. Wariant ten
    jest najlepszy lub remisuje na każdej z~20 instancji i~dorównuje lub pokonuje HLFET/MCP na 19
    z~nich.
    \item \textbf{H1 potwierdzona, ale z~ważnym zastrzeżeniem.} CPM istotnie obniża $C_{max}$
    (25/25 testów), lecz porównanie wykonano przy równej liczbie pokoleń, a~nie przy równym budżecie
    ewaluacji. Przewaga CPM może wynikać z~większego nakładu obliczeniowego na jedno pokolenie.
    \item \textbf{H2 odrzucona w~stronę przeciwną.} Na grafach rzadkich nie dominuje OX, tylko PMX
    (0/14 istotnych testów za OX). Jest to praktyczny wynik potwierdzający przeciwną hipotezę:
    w~tej reprezentacji i~dekoderze PMX jest bezpieczniejszym wyborem operatora od OX.
    \item \textbf{Prawdopodobna przyczyna przewagi PMX.} Naprawa listą zadań gotowych już wymusza
    dopuszczalny porządek topologiczny, więc atut OX w~zachowaniu kolejności zostaje osłabiony.
    Pozycyjne mieszanie PMX skuteczniej eksploruje przestrzeń rozwiązań.
    \item \textbf{Najtrudniejsza instancja: \texttt{robot\_t88}.} Mediana \texttt{pmx\_cpm}=697
    przegrywa tu nieznacznie z~HLFET/MCP=694, choć najlepszy przebieg \texttt{pmx\_cpm}=671 jest lepszy.
    Pokazuje to, że metoda nie dominuje prostych heurystyk w~każdym sensie statystycznym.
    \item \textbf{Wyniki należy czytać przez pryzmat instancji z~zapasem.} Część grafów jest wymuszona:
    wszystkie konfiguracje osiągają LB, więc nie wnoszą sygnału o~przewadze operatorów. Najważniejsze
    różnice pojawiają się na instancjach trudnych, gdzie \texttt{pmx\_cpm} realnie zmniejsza lukę.
    \item \textbf{PMX jest rekomendacją praktyczną dla tego dekodera.} Przewaga PMX nie musi oznaczać
    uniwersalnej dominacji tego operatora, lecz wynika z~połączenia reprezentacji priorytetowej,
    naprawy ready-list i~dekodera listowego. W~tym układzie PMX jest najlepszym wyborem, a~sam PMX
    bez CPM jest zaskakująco konkurencyjny wobec OX z~CPM.
    \item \textbf{Jakość ma koszt obliczeniowy.} CPM daje najlepsze harmonogramy, ale zwiększa liczbę
    ewaluacji i~dominuje czas działania na dużych lub gęstych grafach -- warianty memetyczne są o~rząd
    do dwóch rzędów wielkości droższe od czystego AE (rys.~\ref{fig:cpuscatter}). Najważniejszym
    kolejnym krokiem jest więc porównanie przy równym budżecie ewaluacji oraz tańsza, budżetowana
    wersja lokalnego przeszukiwania.
\end{itemize}

% =====================================================================
\newpage
\begin{thebibliography}{99}
\bibitem{moscato1989} P.~Moscato, \textit{On Evolution, Search, Optimization, Genetic Algorithms and Martial Arts: Towards Memetic Algorithms}, Caltech Concurrent Computation Program, Report C3P 826, 1989.
\bibitem{graham1969} R.~L. Graham, \textit{Bounds on multiprocessing timing anomalies}, SIAM Journal on Applied Mathematics, vol.~17, no.~2, pp.~416--429, 1969.
\bibitem{davis1985} L.~Davis, \textit{Applying adaptive algorithms to epistatic domains}, w: Proc. IJCAI, pp.~162--164, 1985.
\bibitem{goldberg1985} D.~E. Goldberg, R.~Lingle, \textit{Alleles, loci, and the traveling salesman problem}, w: Proc. 1st Int. Conf. on Genetic Algorithms, pp.~154--159, 1985.
\bibitem{kelley1959} J.~E. Kelley, M.~R. Walker, \textit{Critical-path planning and scheduling}, w: Proc. Eastern Joint Computer Conference, pp.~160--173, 1959.
\bibitem{tobita2002} T.~Tobita, H.~Kasahara, \textit{A standard task graph set for fair evaluation of multiprocessor scheduling algorithms}, Journal of Scheduling, vol.~5, no.~5, pp.~379--394, 2002.
\bibitem{adam1974} T.~L. Adam, K.~M. Chandy, J.~R. Dickson, \textit{A comparison of list schedules for parallel processing systems}, Communications of the ACM, vol.~17, no.~12, pp.~685--690, 1974.
\bibitem{wu1990} M.-Y. Wu, D.~D. Gajski, \textit{Hypertool: A programming aid for message-passing systems}, IEEE TPDS, vol.~1, no.~3, pp.~330--343, 1990.
\bibitem{wilcoxon1945} F.~Wilcoxon, \textit{Individual comparisons by ranking methods}, Biometrics Bulletin, vol.~1, no.~6, pp.~80--83, 1945.
\bibitem{deap2012} F.-A. Fortin \textit{et al.}, \textit{DEAP: Evolutionary algorithms made easy}, JMLR, vol.~13, pp.~2171--2175, 2012.
\bibitem{ullman1975} J.~D. Ullman, \textit{NP-complete scheduling problems}, J. Computer and System Sciences, vol.~10, no.~3, pp.~384--393, 1975.
\bibitem{miller1995} B.~L. Miller, D.~E. Goldberg, \textit{Genetic algorithms, tournament selection, and the effects of noise}, Complex Systems, vol.~9, no.~3, pp.~193--212, 1995.
\bibitem{whitley1994} D.~Whitley, V.~S. Gordon, K.~Mathias, \textit{Lamarckian evolution, the Baldwin effect and function optimization}, w: PPSN III, LNCS 866, Springer, pp.~6--15, 1994.
\bibitem{virtanen2020} P.~Virtanen \textit{et al.}, \textit{SciPy 1.0: Fundamental algorithms for scientific computing in Python}, Nature Methods, vol.~17, pp.~261--272, 2020.
\bibitem{mitchell2011} S.~Mitchell, M.~O'Sullivan, I.~Dunning, \textit{PuLP: A linear programming toolkit for Python}, Univ. of Auckland, 2011.
\bibitem{dejong1990} K.~A. De Jong, W.~M. Spears, \textit{An analysis of the interacting roles of population size and crossover in genetic algorithms}, w: PPSN I, LNCS 496, Springer, pp.~38--47, 1990.
\bibitem{grefenstette1986} J.~J. Grefenstette, \textit{Optimization of control parameters for genetic algorithms}, IEEE Trans. Systems, Man, and Cybernetics, vol.~16, no.~1, pp.~122--128, 1986.
\bibitem{omara2010} F.~A. Omara, M.~M. Arafa, \textit{Genetic algorithms for task scheduling problem}, Journal of Parallel and Distributed Computing, vol.~70, no.~1, pp.~13--22, 2010.
\bibitem{hou1994} E.~S.~H. Hou, N.~Ansari, H.~Ren, \textit{A genetic algorithm for multiprocessor scheduling}, IEEE Transactions on Parallel and Distributed Systems, vol.~5, no.~2, pp.~113--120, 1994.
\end{thebibliography}

\end{document}
